\documentclass[../DoAn.tex]{subfiles}
\begin{document}
Chương này tập trung khảo sát hiện trạng của bài toán đặt vé và quản lý sự kiện trên thiết bị di động, từ đó làm rõ các yêu cầu cần thiết cho hệ thống được phát triển trong đồ án. Trên cơ sở phân tích các nền tảng tương tự và đối chiếu với đặc điểm nghiệp vụ của bài toán, chương này xác định những chức năng cốt lõi và các ràng buộc phi chức năng mà hệ thống cần đáp ứng. 

\section{Khảo sát hiện trạng}
\label{section:2.1}
Trong phạm vi đồ án này, khảo sát hiện trạng được thực hiện chủ yếu dựa trên hai nguồn chính là phân tích các ứng dụng tương tự đang tồn tại và đối chiếu với đặc điểm nghiệp vụ của hệ thống \textit{event-booking-app} đang được phát triển. Cách tiếp cận này phù hợp với bối cảnh của một đồ án ứng dụng, khi mục tiêu không chỉ là nhận diện nhu cầu của người dùng cuối mà còn là xác định rõ các khoảng trống chức năng mà hệ thống đề xuất cần giải quyết. Qua việc xem xét luồng nghiệp vụ hiện tại của hệ thống, có thể thấy sản phẩm hướng tới ba nhóm người dùng chính là người tham dự sự kiện, người tổ chức và quản trị viên. Do đó, hiện trạng cần được khảo sát không chỉ ở khía cạnh mua vé, mà còn ở khía cạnh khám phá sự kiện, quản lý thông tin, tương tác giữa các bên và kiểm soát việc tham dự.

\subsection{Đặc thù của bài toán đặt vé và quản lý sự kiện}
Đối với người tham dự, nhu cầu cơ bản không dừng lại ở việc tìm thấy một sự kiện phù hợp, mà còn bao gồm khả năng theo dõi sự kiện yêu thích, nhận thông báo đúng thời điểm, đặt vé nhanh chóng, lưu trữ vé điện tử thuận tiện và được hỗ trợ trong quá trình tham gia sự kiện. Đối với người tổ chức, bài toán lại mở rộng thêm các nhu cầu về tạo sự kiện, cấu hình nhiều hạng vé, quản lý số lượng người tham gia, gửi thông báo và xác nhận check-in tại chỗ. Nếu các khâu này bị tách rời trên nhiều nền tảng khác nhau, trải nghiệm tổng thể sẽ trở nên rời rạc, khó theo dõi và làm tăng chi phí thao tác cho cả người dùng lẫn người tổ chức.

Từ góc nhìn đó, một hệ thống đặt vé sự kiện hiện đại cần thỏa mãn đồng thời ba nhóm chức năng. Nhóm thứ nhất là chức năng khám phá và theo dõi sự kiện, bao gồm tìm kiếm, lọc, xem chi tiết, lưu sự kiện và nhận nhắc lịch. Nhóm thứ hai là chức năng giao dịch và kiểm soát tham dự, bao gồm đặt vé, xác nhận thanh toán, tạo vé điện tử và check-in bằng mã số hoặc mã \gls{QR}. Nhóm thứ ba là chức năng tương tác và hỗ trợ cộng đồng, bao gồm thông báo, lời mời, nhắn tin và các cơ chế kết nối giữa người tham dự với người tổ chức. Đây cũng chính là cơ sở để đánh giá mức độ phù hợp của các ứng dụng tương tự trong phần tiếp theo.

\subsection{Khảo sát một số ứng dụng tương tự}
Trong số các nền tảng hiện có, Meetup và Ticketmaster là hai trường hợp tiêu biểu nhưng đại diện cho hai định hướng khác nhau của bài toán sự kiện. Meetup nhấn mạnh vào việc hình thành cộng đồng, tìm kiếm nhóm theo sở thích và tham gia các hoạt động địa phương hoặc trực tuyến \cite{meetup2026about}. Ngược lại, Ticketmaster tập trung mạnh vào khâu khám phá sự kiện quy mô lớn, mua vé điện tử, quản lý vé trên di động, chuyển vé và bán lại vé trong hệ sinh thái của mình. Việc khảo sát hai nền tảng này giúp nhận diện rõ sự khác biệt giữa hướng tiếp cận thiên về cộng đồng và hướng tiếp cận thiên về thương mại hóa vé điện tử.

\begin{table}[H]
\centering
\renewcommand{\arraystretch}{1.3}
\begin{tabular}{|p{2.7cm}|p{3.5cm}|p{3.8cm}|p{3.8cm}|}
\hline
\textbf{Nền tảng} & \textbf{Điểm mạnh nổi bật} & \textbf{Hạn chế so với bài toán của đồ án} & \textbf{Giá trị tham khảo cho hệ thống} \\
\hline
Meetup \cite{meetup2026about} & Mạnh về khám phá cộng đồng, nhóm sở thích, sự kiện địa phương và kết nối giữa những người có cùng mối quan tâm. & Chức năng đặt vé, quản lý nhiều hạng vé, xác nhận thanh toán và kiểm soát check-in không phải trọng tâm chính; chưa nhấn mạnh mạnh vào quy trình vé điện tử khép kín. & Tham khảo cách tổ chức sự kiện theo cộng đồng, gợi mở nhu cầu kết hợp giữa sự kiện và tương tác xã hội trong ứng dụng. \\
\hline
Ticketmaster & Mạnh về tìm kiếm sự kiện, vé điện tử, quản lý vé trên điện thoại, chuyển vé, bán vé và hỗ trợ người dùng trong ngày diễn ra sự kiện. & Thiên về các sự kiện thương mại quy mô lớn; ít nhấn mạnh chiều tương tác cộng đồng, trò chuyện trực tiếp hoặc nhu cầu tổ chức sự kiện linh hoạt cho nhóm nhỏ. & Tham khảo luồng mua vé, quản lý vé điện tử, truy cập vé trên di động và cơ chế kiểm soát vé thuận tiện. \\
\hline
Giải pháp đề xuất & Hướng tới tích hợp khám phá sự kiện, tạo và quản lý sự kiện, nhiều hạng vé, vé điện tử, thông báo, lời mời, chat và check-in trong cùng một nền tảng di động. & Quy mô triển khai hiện tại còn ở mức thử nghiệm, chưa tối ưu cho lượng người dùng cực lớn hoặc hệ sinh thái phân phối vé quy mô thương mại. & Kết hợp ưu điểm của hướng cộng đồng và hướng quản lý vé điện tử, phù hợp với phạm vi ứng dụng thực tế cỡ vừa. \\
\hline
\end{tabular}
\caption{So sánh một số ứng dụng tương tự với định hướng của giải pháp đề xuất}
\label{table:app-comparison}
\end{table}

Từ bảng so sánh trên có thể thấy, Meetup có ưu thế rõ rệt ở khía cạnh xây dựng cộng đồng và tạo động lực tham gia sự kiện thông qua sở thích chung. Đây là đặc điểm quan trọng vì bài toán sự kiện trong thực tế không chỉ liên quan đến giao dịch mua vé mà còn liên quan tới trải nghiệm kết nối, tương tác và duy trì sự quan tâm của người dùng đối với các hoạt động trong tương lai. Tuy nhiên, nếu chỉ dựa vào hướng tiếp cận như Meetup, hệ thống sẽ thiếu đi những chức năng thiết yếu cho bài toán đặt vé như quản lý nhiều loại vé, xử lý trạng thái thanh toán hay hỗ trợ kiểm soát tham dự bằng vé điện tử.

Ở chiều ngược lại, Ticketmaster thể hiện rõ ưu thế trong quy trình vé điện tử và quản lý việc tham dự trên di động. Việc cho phép người dùng lưu, truy cập, chuyển hoặc sử dụng vé điện tử trực tiếp trên điện thoại là một điểm tham khảo quan trọng đối với Giải pháp đề xuất. Tuy nhiên, mô hình này thiên nhiều hơn về phân phối vé cho các sự kiện lớn, trong khi bài toán của đề tài cần linh hoạt hơn để phục vụ cả các hoạt động cộng đồng, workshop, sự kiện học thuật và các chương trình có quy mô vừa và nhỏ. Điều đó dẫn tới yêu cầu hệ thống không chỉ phải hỗ trợ đặt vé, mà còn phải duy trì khả năng kết nối giữa người dùng và ban tổ chức qua thông báo, lời mời và trao đổi trực tiếp.

\subsection{Nhận xét và các chức năng cần ưu tiên phát triển}
Từ khảo sát hiện trạng, có thể rút ra rằng không có một ứng dụng đơn lẻ nào trong số các nền tảng tham khảo đồng thời bao phủ cân bằng cả ba chiều: khám phá sự kiện, quản lý vé điện tử và tương tác cộng đồng theo phạm vi mà đề tài đặt ra. Vì vậy, Giải pháp đề xuất cần được định hướng như một nền tảng tích hợp, vừa kế thừa khả năng tổ chức và kết nối cộng đồng, vừa hỗ trợ quy trình mua vé và tham dự sự kiện trên thiết bị di động. Trên cơ sở đó, các chức năng cần ưu tiên phát triển gồm quản lý tài khoản người dùng, tìm kiếm và lọc sự kiện, xem chi tiết sự kiện, lưu sự kiện yêu thích, tạo và quản lý sự kiện, quản lý nhiều hạng vé, đặt vé, xác nhận thanh toán, sinh vé điện tử, check-in, gửi thông báo, lời mời và hỗ trợ nhắn tin giữa các bên liên quan. Các chức năng này tạo thành nền tảng cho những phần tổng quan chức năng và đặc tả chức năng sẽ được trình bày ở giai đoạn tiếp theo của báo cáo.

\section{Tổng quan chức năng}
\label{section:2.2}
Từ kết quả khảo sát ở phần \ref{section:2.1}, có thể thấy hệ thống cần hỗ trợ đồng thời nhiều nhóm chức năng khác nhau cho ba tác nhân chính là người dùng, người tổ chức và quản trị viên. Ở mức tổng quan, các chức năng của hệ thống được chia thành bốn nhóm lớn. Nhóm thứ nhất là quản lý tài khoản và hồ sơ người dùng, bao gồm đăng ký, đăng nhập, khôi phục mật khẩu, cập nhật hồ sơ và thiết lập nhận thông báo. Nhóm thứ hai là khám phá và theo dõi sự kiện, bao gồm tìm kiếm, lọc, xem chi tiết, lưu sự kiện yêu thích và xem các sự kiện đã lưu hoặc đã tham gia. Nhóm thứ ba là xử lý giao dịch và tham dự sự kiện, bao gồm đặt vé, xác nhận thanh toán, lưu vé điện tử, hiển thị mã \gls{QR} và kiểm tra tham dự tại sự kiện. Nhóm cuối cùng là hỗ trợ tương tác và vận hành hệ thống, bao gồm tạo sự kiện, quản lý hạng vé, nhắn tin, gửi lời mời, gửi thông báo và duyệt sự kiện ở phía quản trị viên.


\subsection{Biểu đồ use case tổng quát}
\label{subsection:2.2.1}
\begin{figure}[H]
\centering
\includegraphics[width=0.62\linewidth,height=0.78\textheight,keepaspectratio]{Hinhve/UseCaseTongQuat.png}
\caption{Biểu đồ use case tổng quát của hệ thống event-booking-app}
\label{fig:usecase-overview}
\end{figure}

Trong hệ thống có ba tác nhân chính. Tác nhân thứ nhất là người dùng, đại diện cho người tham dự sự kiện. Người dùng thực hiện các chức năng như đăng ký và đăng nhập, tìm kiếm và lọc sự kiện, xem chi tiết sự kiện, lưu sự kiện quan tâm, đặt vé, xác nhận thanh toán, xem vé điện tử, trò chuyện với người tổ chức và nhận thông báo từ hệ thống. Tác nhân thứ hai là người tổ chức, là người có quyền tạo và quản lý sự kiện của mình. Ngoài các chức năng cơ bản như người dùng thông thường, người tổ chức còn có thể tạo sự kiện, cập nhật thông tin sự kiện, quản lý hạng vé, mời người dùng tham gia và quét mã \gls{QR} để check-in cho người tham dự. Tác nhân thứ ba là quản trị viên, chịu trách nhiệm duyệt hoặc từ chối các sự kiện do người tổ chức gửi lên trước khi chúng được hiển thị công khai trong hệ thống.

Người dùng tập trung vào khám phá và tham gia sự kiện, người tổ chức tập trung vào vận hành sự kiện, còn quản trị viên tập trung vào kiểm soát nội dung và trạng thái duyệt. Cách tổ chức này giúp bảo đảm luồng nghiệp vụ rõ ràng, đồng thời hỗ trợ triển khai cơ chế phân quyền tương ứng ở phía backend.

\subsection{Biểu đồ use case phân rã theo nhóm các chức năng}
\label{subsection:2.2.2}
Trong số các use case mức cao, use case “Đặt vé và tham dự sự kiện” là use case quan trọng nhất vì nó kết nối trực tiếp các chức năng khám phá sự kiện, lựa chọn hạng vé, xác nhận thanh toán, phát hành vé điện tử và kiểm soát check-in. Hình~\ref{fig:usecase-booking} trình bày biểu đồ phân rã của use case này.

\begin{figure}[H]
\centering
\includegraphics[width=0.82\linewidth,height=0.78\textheight,keepaspectratio]{Hinhve/UseCasePhanRa.png}
\caption{Biểu đồ phân rã theo nhóm các use case của hệ thống}
\label{fig:usecase-booking}
\end{figure}
Theo Hình~\ref{fig:usecase-booking}, các use case của hệ thống được nhóm thành những cụm chức năng chính thay vì tách rời từng thao tác đơn lẻ. Nhóm “Quản lý người dùng” bao gồm các chức năng đăng ký, đăng nhập và cập nhật hồ sơ. Nhóm “Khám phá và theo dõi sự kiện” bao gồm tìm kiếm, lọc, xem chi tiết và lưu sự kiện. Nhóm “Đặt vé và quản lý vé” tập trung vào lựa chọn loại vé, thanh toán và xem vé đã đặt. Nhóm “Tương tác và thông báo” phản ánh các chức năng chat và nhận thông báo. Với phía người tổ chức, nhóm “Quản lý sự kiện” bao gồm tạo sự kiện, quản lý sự kiện, quản lý hạng vé, mời người dùng và check-in vé. Cuối cùng, phía quản trị viên đảm nhiệm nhóm chức năng kiểm duyệt hệ thống thông qua việc duyệt hoặc từ chối sự kiện. Cách phân nhóm này giúp làm rõ cấu trúc nghiệp vụ tổng thể của hệ thống và là cơ sở thuận lợi để lựa chọn các use case quan trọng cần đặc tả ở phần sau.

\subsection{Quy trình nghiệp vụ}
\label{subsection:2.2.3}
\begin{figure}[H]
\centering
\includegraphics[width=0.75\linewidth,height=0.82\textheight,keepaspectratio]{Hinhve/QuyTrinhNghiepVu.png}
\caption{Biểu đồ hoạt động của quy trình nghiệp vụ tham gia sự kiện}
\label{fig:activity-booking}
\end{figure}

Quy trình nghiệp vụ tiêu biểu của hệ thống là quy trình tham gia sự kiện, được mô tả trong Hình~\ref{fig:activity-booking}. Đây là một quy trình kết hợp nhiều use case khác nhau, bắt đầu từ lúc người dùng tìm kiếm sự kiện, xem chi tiết, có thể lưu lại để theo dõi, sau đó lựa chọn loại vé và thực hiện thanh toán. Sau khi thanh toán thành công, hệ thống phát hành vé điện tử và hỗ trợ gửi thông báo xác nhận. Trước thời điểm diễn ra chương trình, người dùng có thể nhận thông báo nhắc lịch và nếu cần có thể trao đổi thêm với người tổ chức qua chức năng chat. Khi đến địa điểm tổ chức, người dùng xuất trình vé điện tử có mã \gls{QR}; người tổ chức quét mã để hệ thống xác thực và cập nhật trạng thái check-in. Như vậy, quy trình này không chỉ giới hạn ở một thao tác mua vé, mà phản ánh đầy đủ một chuỗi hoạt động liên kết từ khám phá, giao dịch, theo dõi cho đến tham dự thực tế.
\section{Đặc tả chức năng}
\label{section:2.3}
Trên cơ sở biểu đồ tổng quan và quy trình nghiệp vụ ở phần \ref{section:2.2}, đồ án lựa chọn bốn use case quan trọng để đặc tả chi tiết. Đây là các use case phản ánh rõ nhất những nghiệp vụ cốt lõi của hệ thống, bao gồm phía người tổ chức, người dùng và quá trình tương tác giữa các tác nhân trong suốt vòng đời của một sự kiện. Bốn use case được lựa chọn gồm: tạo sự kiện, đặt vé và xác nhận thanh toán, gửi lời mời và thông báo, cùng với check-in vé bằng mã \gls{QR}.

\subsection{Đặc tả use case Tạo sự kiện}
\textbf{Tên use case:} Tạo sự kiện.

\textbf{Tiền điều kiện:} Người tổ chức đã đăng nhập thành công và có quyền sử dụng chức năng tạo sự kiện.

\textbf{Hậu điều kiện:} Một sự kiện mới được lưu trong hệ thống ở trạng thái chờ duyệt, đồng thời các hạng vé ban đầu của sự kiện cũng được khởi tạo.

\textbf{Luồng sự kiện chính:}
Người tổ chức truy cập chức năng tạo sự kiện và nhập các thông tin cần thiết như tiêu đề, mô tả, danh mục, thời gian, địa điểm, hình ảnh và danh sách hạng vé. Hệ thống kiểm tra tính hợp lệ của dữ liệu đầu vào, bao gồm các trường bắt buộc cũng như thông tin của từng hạng vé như tên vé, giá vé và số lượng vé. Khi toàn bộ thông tin hợp lệ, hệ thống tạo bản ghi sự kiện mới, gán người tạo là người tổ chức hiện tại, thiết lập trạng thái duyệt ban đầu là chờ duyệt và lưu các hạng vé tương ứng.

\textbf{Luồng phát sinh:}
Nếu người tổ chức bỏ trống các trường bắt buộc như tiêu đề, danh mục, ngày hoặc địa điểm, hệ thống từ chối lưu và yêu cầu bổ sung thông tin. Nếu danh sách hạng vé không hợp lệ, chẳng hạn không có hạng vé nào, giá vé âm hoặc quota không hợp lệ, hệ thống thông báo lỗi và không tạo sự kiện. Nếu dữ liệu gửi lên sai định dạng, hệ thống trả về thông báo lỗi tương ứng.

\subsection{Đặc tả use case Đặt vé và xác nhận thanh toán}
\textbf{Tên use case:} Đặt vé và xác nhận thanh toán.

\textbf{Tiền điều kiện:} Người dùng đã đăng nhập; sự kiện tồn tại trong hệ thống; hạng vé được chọn hợp lệ và còn đủ số lượng theo quota.

\textbf{Hậu điều kiện:} Giao dịch thanh toán được ghi nhận; vé được phát hành ở dạng điện tử; mã \gls{QR} của vé được sinh ra để phục vụ việc tham dự sự kiện.

\textbf{Luồng sự kiện chính:}
Người dùng chọn một sự kiện mong muốn, sau đó lựa chọn hạng vé và số lượng vé cần mua. Hệ thống kiểm tra tính hợp lệ của lựa chọn, tính tổng giá trị giao dịch và khởi tạo các bản ghi vé ở trạng thái chờ thanh toán cùng với bản ghi thanh toán tương ứng. Người dùng tiếp tục xác nhận thanh toán. Khi thanh toán thành công, hệ thống cập nhật trạng thái giao dịch, đánh dấu vé là đã thanh toán, tăng số lượng vé đã bán của hạng vé tương ứng, đồng thời sinh dữ liệu \gls{QR} cho từng vé điện tử. Sau đó người dùng có thể xem lại vé đã đặt trong ứng dụng.

\textbf{Luồng phát sinh:}
Nếu hạng vé không tồn tại hoặc không thuộc sự kiện đang chọn, hệ thống dừng thao tác và thông báo lỗi. Nếu số lượng vé yêu cầu vượt quá quota còn lại, hệ thống từ chối đặt vé. Nếu tổng giá tiền gửi lên không khớp với cấu hình giá vé của hệ thống, hệ thống trả về lỗi xác thực giao dịch. Nếu thanh toán không thành công, các vé tương ứng được giữ ở trạng thái thất bại hoặc không được xác nhận hoàn tất.

\subsection{Đặc tả use case Gửi lời mời và thông báo}
\textbf{Tên use case:} Gửi lời mời và thông báo.

\textbf{Tiền điều kiện:} Tác nhân thực hiện đã đăng nhập; sự kiện cần gửi lời mời tồn tại; danh sách người nhận hợp lệ.

\textbf{Hậu điều kiện:} Các thông báo hoặc lời mời được tạo trong hệ thống; nếu phù hợp, thông báo đẩy được gửi tới thiết bị của người nhận.

\textbf{Luồng sự kiện chính:}
Người dùng hoặc người tổ chức chọn một sự kiện và xác định danh sách người nhận muốn mời tham gia. Hệ thống kiểm tra sự tồn tại của sự kiện, sau đó tạo các bản ghi thông báo cho từng người nhận với nội dung lời mời tương ứng. Nếu người nhận đã bật thông báo và đã đăng ký thiết bị, hệ thống tiếp tục gửi thông báo đẩy tới thiết bị di động của họ. Ngoài lời mời sự kiện, cùng cơ chế này còn được dùng để ghi nhận và gửi các thông báo hệ thống như xác nhận thanh toán, tin nhắn mới hoặc nhắc lịch.

\textbf{Luồng phát sinh:}
Nếu sự kiện không tồn tại, hệ thống hủy thao tác và thông báo lỗi. Nếu danh sách người nhận rỗng hoặc không hợp lệ, hệ thống không tạo lời mời. Nếu người nhận đã tắt thông báo hoặc chưa đăng ký thiết bị nhận thông báo, hệ thống vẫn có thể lưu bản ghi thông báo nội bộ nhưng không gửi thông báo đẩy ra bên ngoài.

\subsection{Đặc tả use case Check-in vé bằng mã QR}
\textbf{Tên use case:} Check-in vé bằng mã \gls{QR}.

\textbf{Tiền điều kiện:} Người tổ chức đã đăng nhập; là chủ sở hữu của sự kiện tương ứng; người tham dự đã có vé hợp lệ ở trạng thái thanh toán thành công.

\textbf{Hậu điều kiện:} Vé được cập nhật trạng thái đã check-in nếu hợp lệ; ngược lại hệ thống từ chối check-in và giữ nguyên trạng thái vé.

\textbf{Luồng sự kiện chính:}
Người tổ chức mở chức năng check-in và quét mã \gls{QR} từ vé điện tử của người tham dự. Dữ liệu quét được gửi tới hệ thống để phân tích. Hệ thống kiểm tra mã vé có đúng định dạng hay không, vé có thuộc sự kiện hiện tại hay không, trạng thái thanh toán của vé đã hợp lệ hay chưa và vé đã từng được sử dụng trước đó chưa. Nếu tất cả điều kiện đều thỏa mãn, hệ thống đánh dấu vé là đã check-in, ghi nhận thời điểm check-in và người thực hiện xác nhận. Sau đó hệ thống trả về kết quả check-in thành công cho người tổ chức.

\textbf{Luồng phát sinh:}
Nếu mã \gls{QR} sai định dạng, hệ thống từ chối xử lý và thông báo mã không hợp lệ. Nếu vé không thuộc sự kiện hiện tại, hệ thống thông báo sai sự kiện. Nếu vé chưa được thanh toán, hệ thống từ chối check-in. Nếu vé đã được sử dụng trước đó, hệ thống trả về trạng thái đã check-in và không cập nhật lại dữ liệu.

\section{Yêu cầu phi chức năng}
\label{section:2.4}
Ngoài các chức năng nghiệp vụ, hệ thống còn phải đáp ứng một số yêu cầu phi chức năng để bảo đảm khả năng vận hành ổn định, an toàn và thuận tiện cho người sử dụng. Các yêu cầu này được xác định dựa trên đặc điểm của ứng dụng di động có tích hợp backend, xử lý dữ liệu sự kiện, vé điện tử và tương tác gần thời gian thực.

\subsection{Yêu cầu về hiệu năng}
Hệ thống cần bảo đảm thời gian phản hồi hợp lý đối với các chức năng thường xuyên được sử dụng như đăng nhập, tải danh sách sự kiện, xem chi tiết sự kiện, gửi tin nhắn và truy xuất vé điện tử. Đối với người dùng di động, trải nghiệm bị ảnh hưởng mạnh nếu thời gian tải dữ liệu quá lâu hoặc thao tác xác nhận thanh toán, mở vé và check-in diễn ra chậm. Vì vậy, ứng dụng cần tổ chức API gọn nhẹ, hạn chế truy vấn dư thừa và tối ưu số lượng dữ liệu trả về ở các màn hình chính.

\subsection{Yêu cầu về độ tin cậy và tính nhất quán dữ liệu}
Do hệ thống xử lý nhiều thực thể liên quan như người dùng, sự kiện, vé, thanh toán và thông báo, dữ liệu cần được cập nhật nhất quán giữa ứng dụng khách và máy chủ. Trạng thái vé, số lượng vé đã bán, thông tin thanh toán và trạng thái check-in phải phản ánh đúng thực tế tại thời điểm sử dụng. Hệ thống cũng cần giảm thiểu lỗi phát sinh khi nhiều người dùng đồng thời thao tác trên cùng một sự kiện, đặc biệt trong các trường hợp đặt vé hoặc xác nhận tham dự.

\subsection{Yêu cầu về bảo mật}
Hệ thống cần bảo vệ thông tin tài khoản và dữ liệu nghiệp vụ thông qua cơ chế xác thực, phân quyền và kiểm soát truy cập phù hợp. Các API liên quan đến tài khoản, vé, thanh toán, tin nhắn và quản trị cần chỉ cho phép truy cập khi người dùng đã được xác thực hợp lệ. Ngoài ra, các dữ liệu nhạy cảm như mật khẩu, token và thông tin thao tác của người dùng phải được quản lý an toàn, tránh lộ lọt hoặc bị sử dụng trái phép.

\subsection{Yêu cầu về khả năng sử dụng}
Ứng dụng hướng tới người dùng cuối thao tác chủ yếu trên điện thoại, vì vậy giao diện cần rõ ràng, dễ hiểu và hạn chế các bước xử lý phức tạp. Các chức năng quan trọng như tìm sự kiện, đặt vé, xem vé đã mua, nhận thông báo và gửi tin nhắn cần được bố trí trực quan. Hệ thống cũng cần hỗ trợ trải nghiệm liên tục giữa các màn hình, giúp người dùng không bị đứt quãng khi chuyển từ khám phá sự kiện sang đặt vé hoặc từ thông báo sang xem chi tiết nội dung liên quan.

\subsection{Yêu cầu về khả năng bảo trì và mở rộng}
Do đây là hệ thống có nhiều nhóm chức năng, mã nguồn cần được tổ chức theo cấu trúc rõ ràng để thuận lợi cho việc bảo trì và mở rộng sau này. Các thành phần như xác thực, sự kiện, vé, thông báo, nhắn tin và nhắc lịch nên được tách thành các mô-đun tương đối độc lập, giúp dễ sửa lỗi, bổ sung tính năng và kiểm thử. Yêu cầu này đặc biệt quan trọng vì hệ thống có thể được mở rộng thêm các cổng thanh toán thực tế, công cụ quản trị nâng cao hoặc cơ chế gợi ý cá nhân hóa trong các phiên bản tiếp theo.

\begin{samepage}
\subsection{Yêu cầu kỹ thuật triển khai}
Hệ thống cần có khả năng hoạt động trên môi trường di động phổ biến và kết nối được với backend thông qua mạng Internet thông thường. Về phía lưu trữ, cơ sở dữ liệu cần đủ linh hoạt để quản lý các thực thể thay đổi theo nghiệp vụ như sự kiện, hạng vé, thông báo và phòng chat. Về phía tích hợp, hệ thống cũng cần hỗ trợ các dịch vụ bổ sung như thông báo đẩy, sinh mã \gls{QR} và xử lý nội dung hỗ trợ bằng \gls{API}, từ đó tạo nền tảng kỹ thuật phù hợp cho việc hiện thực hóa sản phẩm.

%%%%%%%%%%%%%%%%%%%%%%%%%%%%%%%%%%%

Từ kết quả khảo sát trong chương này có thể thấy bài toán đặt vé và quản lý sự kiện trên thiết bị di động đòi hỏi một hệ thống tích hợp nhiều chức năng hơn so với các nền tảng chỉ thiên về cộng đồng hoặc bán vé. Việc phân tích các ứng dụng tương tự đã giúp xác định rõ những chức năng cốt lõi cần ưu tiên, làm rõ các yêu cầu phi chức năng mà hệ thống phải đáp ứng để vận hành ổn định trong thực tế.
\end{samepage}

\end{document}
